Ein Klotz von \mKO\ hängt in Ruhe an einer Feder, die dabei um \xoO\ gedehnt ist. Man hebt ihn an, bis die Feder gerade entspannt ist, und lässt ihn dort aus der Ruhe los. Luftwiderstand und die Masse der Feder sind vernachlässigbar. Rechne mit \gO.

\begin{abcliste}
\abc
Wie gross ist die Federkonstante?

\abc
Wie weit sinkt der Klotz, bis er umkehrt? Stelle eine Energiebilanz auf.

\abc
Wo ist der Klotz am schnellsten, und wie schnell ist er dort?

\abc
Was gilt im tiefsten Punkt für die Geschwindigkeit und für die Kräfte? Warum bleibt der Klotz dort nicht stehen?
\end{abcliste}
